\section{関門の判定と下位指標の定義}\label{app:gates}

\paragraph{自己位置の追跡（G1）}
運動が見えるステップで，$h^1$ から A-1 の位置を線形回帰で読み，決定係数が 0.80 以上であることを前提とした．

\paragraph{価値と行動・色の結びつき（G2）}
運動が隠れたステップで，次の三つを測った．
M1(c) は，A-1 が色 $X$ の象限の方へ（進む向きと $X$ のランドマークの方向の角度差が 45° 以内）動いたステップの予測誤差から，それ以外へ動いたステップの予測誤差を引いた値 $\Delta_X$ を求め，$\Delta_{\text{赤}}-\frac13(\Delta_{\text{緑}}+\Delta_{\text{青}}+\Delta_{\text{シアン}})$ としたものである．
自分が赤の方へ動くと見込んでいれば負になる．
M1(d) は，A-1 が実際には $X$ の方へ動かなかったステップで，予測が実際の次の視界より「$X$ の方へ一歩動いた場合の視界」（シミュレータで描き直したもの）に近いかを色ごとに測り（寄り $S_X$），赤に特有の寄りを取り出したものである．
赤の方へ動くと見込んでいれば正になる．
% 【要確認】M1(d) の c_cf が S_赤 − (ほか3色の平均) か，S_赤 そのものか（v11 ログ §7.9）。第2弾で確かめて式を書く。
色の基準は，赤と緑を入れ替えた配置のエピソードで，寄りが新しい赤の方へ向かうか（$S_{\text{赤}}-S_{\text{緑}}$）を測るものである．
休止のステップは除き，各色の比較からはその色のランドマークから距離2以内のステップを除いた．
予備的な試み（快・不快を入力として一から学習したモデル）では，これらの大きさは予測誤差の 5〜8\% ほどで，赤・緑の画素を除いても寄りの 80〜90\% が残った．本文の条件のモデルでは，この確認は行っていない．
% 出典：v11 まとめ §2.1〜§2.2，v11 ログ §7.6・§7.9，v11 まとめ §5 の 5。

\paragraph{自分側・他者側の積み上がり（G3・G5a）と符号応答（G5b）}
G3 と G5b は本文3.4節で定義した．
G5a は，A-2 の本当の快・不快を $t=20$ の値に保って15ステップ進めたときの，予測の中の A-2 の傾き $g_{\mathrm{other}}$ である．

\paragraph{判定}
各指標の判定を表\ref{tab:gates}に示す．

\begin{table}[tbp]
  \centering
  \small
  \caption{段階Iの関門の判定}
  \label{tab:gates}
  % paper/tables/gates.tex -- generated by analyze/paper/tables_paper.py; sources in paper/tables/README.md
\begin{tabular}{@{}llllllll@{}}
\toprule
指標 & 条件 & 組1 & 組2 & 組3 & 組4 & 組5 & 組6 \\
\midrule
自己位置の追跡（前提） & 経験50\% & 5/5 & 5/5 & 5/5 & 5/5 & 5/5 & 5/5 \\
 & 経験80\% & 0/5 & 0/5 & 0/5 & 0/5 & 0/5 & 0/5 \\
 & 非共有 & 0/5 & 0/5 & 0/5 & --- & --- & --- \\
\midrule
価値と行動・色の結びつき & 経験50\% & 5/5 & 5/5 & 5/5 & 5/5 & 5/5 & 5/5 \\
 & 経験80\% & 5/5 & 5/5 & 5/5 & 5/5 & 5/5 & 5/5 \\
 & 非共有 & \multicolumn{3}{l}{3組で合格} & --- & --- & --- \\
\midrule
自己側の積み上がり & 経験50\% & 5/5 & 5/5 & 5/5 & 5/5 & 5/5 & 5/5 \\
 & 経験80\% & 5/5 & 5/5 & 5/5 & 5/5 & 5/5 & 5/5 \\
 & 非共有 & 5/5 & 5/5 & 5/5 & --- & --- & --- \\
\midrule
他者側の積み上がり（G5a） & 経験50\% & 5/5 & 4/5 & 3/5 & 1/5 & 5/5 & 0/5 \\
 & 経験80\% & 5/5 & 0/5 & 1/5 & 0/5 & 0/5 & 0/5 \\
 & 非共有 & 0/5 & 0/5 & 0/5 & --- & --- & --- \\
\midrule
他者側の符号応答（G5b） & 経験50\% & 5/5 & 2/5 & 3/5 & 0/5 & 0/5 & 5/5 \\
 & 経験80\% & 5/5 & 5/5 & 5/5 & 0/5 & 3/5 & 5/5 \\
 & 非共有 & 0/5 & 0/5 & 0/5 & --- & --- & --- \\
\bottomrule
\end{tabular}

  % 表A1：関門の判定（自己位置の追跡，G2，G3，G5a，G5b）を，経験50%・経験80%（組1〜6），非共有条件（組1〜3）で。
  % 各欄は合格の点の数（5点中）と組の数。出典：v13 まとめ §3 表1・§5 表4・§10。
\end{table}
